\appendix

\section{Proof details}
\label{app:proofs}

This appendix collects the typed arguments supporting the concise statements
in the main text. All scalar sums are XOR reductions and all scalar products
are conjunctions unless another algebra is named explicitly.

\subsection{Representation and selection}

\begin{proposition}[Unique binary representation]
The map from a binary Boolean operator $\Theta$ to $\CM{\Theta}$ is a
bijection between functions $\Bool^2\rightarrow\Bool$ and matrices in
$\Bool^{2\times2}$ once the row, column, and truth-state order is fixed.
\end{proposition}

\begin{proof}
A binary Boolean function is determined by its four values on
$(1,1),(1,0),(0,1),(0,0)$. These are exactly the four entries of
$\CM{\Theta}$. Conversely, any four-bit matrix assigns one output to each
binary input pair and therefore defines one such function. Both sets contain
$2^4=16$ elements.
\end{proof}

For completeness, expand the contraction as
\begin{equation}
  E_{\CM{\Theta}}(x,y)
  =\mathop{\XOR}_{r,s\in\{1,0\}}
    \chi_x(r)\wedge\Theta_{rs}\wedge\chi_y(s),
  \qquad
  \chi_z(q):=\begin{cases}1,&q=z,\\0,&q\ne z.\end{cases}
\end{equation}
The two one-hot states make
$\chi_x(r)\wedge\chi_y(s)$ equal to one only at $(r,s)=(x,y)$.
This proves selection without a cancellation or distributivity assumption
beyond the declared Boolean algebra. It also explains why OR--AND selection
would return the same scalar on one-hot states, although OR does not supply the
same coefficient-space algebra for arbitrary vectors.

\subsection{Coefficient linearity and basis}

For fixed $x,y$, distribute conjunction over XOR:
\begin{align}
  E_{A\XOR B}(x,y)
  &=\mathop{\XOR}_{r,s}
    \chi_x(r)\wedge(A_{rs}\XOR B_{rs})\wedge\chi_y(s) \\
  &=E_A(x,y)\XOR E_B(x,y).
\end{align}
Thus evaluation is a coordinate projection and hence a linear functional on
$\GFtwo^{2\times2}$. Let $E_{rs}$ be the matrix with a one only at position
$(r,s)$. Evaluating
\begin{equation}
  \CM{\Theta}=\mathop{\XOR}_{r,s}\Theta_{rs}E_{rs}
\end{equation}
at any position selects $\Theta_{rs}$, proving the basis decomposition and
coefficient extraction. This proof concerns the coefficient matrix. The
affine truth-state encoding prevents it from implying that every represented
Boolean function is linear in its logical inputs.

\subsection{Signed operand transformations}

Let $P=\begin{bmatrix}0&1 \\ 1&0\end{bmatrix}$. The row-$x$, column-$y$
entry of $P\CM{\Theta}$ is the row-$\neg x$, column-$y$ entry of
$\CM{\Theta}$ and is therefore $(\neg x)\mathbin{\Theta}y$. Similarly,
right multiplication by $P$ exchanges the column assignment and
transposition exchanges the two input assignments. With
$\Theta^\tau(x,y):=\Theta(y,x)$, hence
\begin{align}
  \CM{\Theta^\tau} &= \CM{\Theta}^{\mathsf T}, &
  \CM{\Theta(\neg X,Y)} &= P\CM{\Theta}, \\
  \CM{\Theta(X,\neg Y)} &= \CM{\Theta}P, &
  \CM{\Theta(\neg X,\neg Y)} &= P\CM{\Theta}P.
\end{align}
Negating every output flips every bit, which is entrywise XOR with the
all-ones matrix. Products with $P$ are XOR--AND products, but only one entry
survives in each output position, so they can equivalently be implemented as
direct permutations.

For XNOR, $\CM{\Leftrightarrow}$ is the identity matrix. Under the explicit
clockwise convention
$\operatorname{rot}_{90^\circ}\!\begin{bmatrix}a&b\\c&d\end{bmatrix}
=\begin{bmatrix}c&a\\d&b\end{bmatrix}$, its ones move to the antidiagonal,
giving $\CM{\XOR}$. (For this particular symmetric matrix, counterclockwise
rotation gives the same result.) The two supports are disjoint and cover all four
positions, so their entrywise XOR is the all-ones $\CM{\impax}$.

\subsection{Operand-aligned fusion}

Let $\Phi$ be any binary Boolean operator and define
\begin{equation}
  F_{rs}:=(\Theta_1)_{rs}\mathbin{\Phi}(\Theta_2)_{rs}.
\end{equation}
When both inner operators use the same operand frame, selection gives
\begin{align}
 &(x\mathbin{\Theta_1}y)\mathbin{\Phi}
   (x\mathbin{\Theta_2}y) \\
 &\qquad=(\Theta_1)_{xy}\mathbin{\Phi}(\Theta_2)_{xy}
 =F_{xy}
 =\statebra{x}F\stateket{y}.
\end{align}
No general distributive law for $\Phi$ over XOR is used. If the original
operators differ by signed operand order, the preceding transformations first
translate each one into the target frame; applying the same argument after
that translation proves fusion after alignment.

For pure structural compilation, proceed by induction on the derivation of
$\Gamma\vdash e\Downarrow^{\mathsf S}\operatorname{Pair}(R,C,t)$. A primitive
signed-literal pair is sound by the transformations above. For the unary step,
complementing all four token bits implements logical negation. For the binary
step, the induction hypotheses give two sound child pairs in the same typed
frame, and the preceding entry-selection calculation proves their entrywise
outer-connective fusion. Every recursive call visits a strict AST child,
proving termination. This establishes structural soundness without using
retabulation.

For exact pair retabulation, evaluate the subtree independently at
$(1,1),(1,0),(0,1),(0,0)$ and store the four observed bits in that order.
Selection at any assignment returns the stored value by construction. The
whole-root retabulation is therefore sound. For a hybrid derivation, induct on
the judgment in \eqref{eq:pair-judgment}. The new base case is an exactly
retabulated subtree. Every negation or fusion above that base preserves the
invariant by the same unary and binary arguments used for pure structural
compilation. Provenance changes the reported construction class, not the
invariant. The top-level compiler consequently returns only a pure structural
pair, a whole-root retabulation, a hybrid pair assembled from sound children,
or the ordinary IR result. The ordinary IR retains responsibility for every
expression not admitted by a pair judgment.

\subsection{Formula-valued LMs}

Let $B_{ij}(X,Y)=\otbktwo{X_i}{Y_j}$. The upper-left cell of
\begin{equation}
  [\mathcal{M}_{X\Theta Y}]
  =\mathop{\XOR}_{i,j\in\{1,2\}}c_{ij}B_{ij}(X,Y),
  \qquad c_{ij}:=\Theta_{b_i b_j},\quad b_1=1,\ b_2=0,
\end{equation}
is the disjoint-minterm expansion
\begin{equation}
  c_{11}XY
  \XOR c_{12}X\neg Y
  \XOR c_{21}\neg X Y
  \XOR c_{22}\neg X\neg Y.
\end{equation}
Under the declared position-to-assignment map, this is
$X\mathbin{\Theta}Y$. Exchanging the $Y$ components gives
$X\mathbin{\Theta}\neg Y$; exchanging the $X$ components gives
$\neg X\mathbin{\Theta}Y$; and exchanging both gives the lower-right
formula. This proves the polarity-orbit form in \eqref{eq:lm-matrix}.

If $v(X)=x$ and $v(Y)=y$, applying $v$ to those four formulas either preserves
or exchanges the two row and column assignments. It follows that
\begin{equation}
  v([\mathcal{M}_{X\Theta Y}])
  =P^{1-x}\CM{\Theta}P^{1-y}.
\end{equation}
The positive case $x=y=1$ yields $\CM{\Theta}$.

For formulas $A,B$,
\begin{equation}
  \statebra{A}\stateket{B}
  =(A\wedge B)\XOR(\neg A\wedge\neg B)
  =(A\Leftrightarrow B),
\end{equation}
because the two conjunctions are disjoint. Substituting the LM basis expansion
between $\statebra{\mathcal L_1}$ and $\stateket{\mathcal L_2}$ and using this
identity on both sides of each outer product gives
\begin{equation}
  \mathop{\XOR}_{i,j}
    c_{ij}\wedge(\mathcal L_1\Leftrightarrow X_i)
    \wedge(Y_j\Leftrightarrow\mathcal L_2).
\end{equation}
For each valuation of $A,B,X,Y$, exactly one $i$ satisfies
$A\Leftrightarrow X_i$ and exactly one $j$ satisfies
$Y_j\Leftrightarrow B$. If these positions are $i^*,j^*$, the XOR therefore
selects $c_{i^*j^*}=\Theta_{b_{i^*}b_{j^*}}$. By the definitions of the
polarity orbit and $b_i$ under this valuation,
$b_{i^*}=v(A\Leftrightarrow X)$ and
$b_{j^*}=v(Y\Leftrightarrow B)$. Hence the selected coefficient equals
$v((A\Leftrightarrow X)\mathbin{\Theta}(Y\Leftrightarrow B))$. This holds for
every valuation, proving \eqref{eq:logical-pairing-identity}.

Finally, a Boolean valuation is a homomorphism for negation, conjunction, and
XOR. Moving it through this finite expression term by term proves
\eqref{eq:valuation-commutes}.

\subsection{ANF conversion is a different map}

For $n$ variables, identify a valuation with its support
$S\subseteq\{1,\ldots,n\}$ and let $t_S$ be the corresponding truth value.
If
\begin{equation}
  f(x_1,\ldots,x_n)
  =\mathop{\XOR}_{T\subseteq\{1,\ldots,n\}}
    \left(a_T\wedge\bigwedge_{i\in T}x_i\right),
\end{equation}
then evaluation on support $S$ gives
\begin{equation}
  t_S=\mathop{\XOR}_{T\subseteq S}a_T.
\end{equation}
Boolean M\"obius inversion gives the converse relation
\begin{equation}
  a_T=\mathop{\XOR}_{S\subseteq T}t_S.
\end{equation}
Thus ANF coefficients arise by an invertible coordinate transform of the truth
vector. By contrast, positive valuation maps the formula-valued polarity orbit
directly to a numeric CM. The two maps have different domains and purposes
even though both ultimately describe the same Boolean function.

\subsection{Higher-arity size}

An arbitrary function of $n$ Boolean variables has one output for each of
$2^n$ valuations. Its explicit operator therefore has $2^n$ bits. The same
data can be placed in an order-$n$ tensor with side length two, a
length-$2^n$ vector, or, after splitting $2m$ variables into equal row and
column groups, a $2^m\times2^m$ matrix. These are layout choices with the same
entry count, not compression results.

\section{Finite verification scope}
\label{app:finite-checks}

The accompanying checker exhaustively verifies the 16 binary tokens, 16,384
same-frame fusion cases, 512 signed alignments, 1,048,576 signed
alignment/fusion cases, the XOR/XNOR/Impax identities, LM positive and general
valuations, 256 LM coefficient extractions, 16,384 LM valuation/fusion cases,
and 64 two-variable ANF round trips. Finite enumeration catches convention and
implementation errors but is not used as a proof of an unbounded theorem or a
runtime bound.
